\section{Questão 05: Geração Aumentada por Recuperação (RAG)}
\label{sec:q5}

\subsection{Objetivo e arquitetura}

A Questão 05 constrói uma aplicação de RAG do tipo \textit{Standard}
(Seção~\ref{sec:fund-rag}) sobre o dataset docentesDC, com o objetivo de fornecer
respostas ancoradas nos materiais acadêmicos reais dos docentes do Departamento de
Computação, reduzindo as alucinações de um LLM genérico consultado sem contexto. O
pipeline, implementado com a biblioteca LangChain, tem cinco etapas: carga e análise
exploratória dos dados; consulta \textit{baseline} ao LLM sem contexto (para
evidenciar a alucinação); engenharia de RAG (fragmentação, \textit{embeddings} e
indexação); comparação prática; e avaliação por \textit{benchmark}. A
Tabela~\ref{tab:q5-config} resume os componentes.

\begin{table}[ht]
  \centering
  \caption{Componentes do pipeline de RAG (Questão 05).}
  \label{tab:q5-config}
  \footnotesize
  \begin{tabularx}{0.92\textwidth}{@{}l X@{}}
    \toprule
    \textbf{Item} & \textbf{Valor} \\
    \midrule
    Tipo de RAG & \textit{Standard} (recuperação + geração) \\
    LLM gerador & \texttt{gpt-4o-mini} \\
    \textit{Embeddings} & \texttt{text-embedding-3-small} (OpenAI) \\
    Banco vetorial & ChromaDB persistente \\
    Fragmentação & \texttt{RecursiveCharacterTextSplitter}, \textit{chunk} de 1000 caracteres, \textit{overlap} de 200 \\
    Filtragem & filtro por metadados (nome do docente) para elevar a precisão \\
    \textit{Prompt} & restritivo: responder apenas com base no contexto recuperado \\
    \bottomrule
  \end{tabularx}
\end{table}

\subsection{Avaliação com Ragas}

A qualidade foi medida sobre um \textit{benchmark} (\textit{golden dataset}) de 30
perguntas com respostas de referência, usando o \textit{framework}
Ragas~\cite{es2024ragas} com um LLM como juiz. A Figura~\ref{fig:q5-ragas} apresenta
as três métricas: fidelidade 0,674, relevância da resposta 0,561 e precisão do
contexto 0,430, todas abaixo do patamar ideal (0,85 para fidelidade).

\begin{figure}[ht]
  \centering
  \includegraphics[width=0.8\textwidth]{q5_ragas.png}
  \caption{Pontuações médias do Ragas no \textit{benchmark} de 30 perguntas. A linha
  tracejada indica o patamar de referência de fidelidade (0,85).}
  \label{fig:q5-ragas}
\end{figure}

\subsection{Diagnóstico das pontuações}

Uma análise linha a linha revelou que as notas baixas decorrem menos de falha do
pipeline e mais de três fatores. Primeiro, há \textbf{desalinhamento entre o
\textit{benchmark} e o corpus}: várias perguntas assumem informações inexistentes no
docentesDC (por exemplo, conceitos de compiladores atribuídos a um docente cujos 275
documentos tratam apenas de gestão de projetos). Quando a informação não existe, o
\textit{prompt} restritivo corretamente faz o modelo responder ``não encontrei essa
informação no contexto fornecido'', mas o Ragas atribui relevância e precisão zero a
essa resposta honesta, derrubando as médias. Segundo, há uma \textbf{discrepância de
metadados na avaliação}: o nome do docente é injetado no \textit{prompt} durante a
execução, mas apenas o texto cru é salvo para o Ragas, levando o juiz a interpretar
como alucinação uma resposta que seguiu fielmente o contexto. Terceiro, ocorre
\textbf{colisão semântica}: em perguntas gerais, planos de ensino (que contêm termos
genéricos e o nome do docente) superam em similaridade os slides detalhados, que
acabam omitidos do contexto.

\subsection{Contribuição do RAG}

Apesar das pontuações limitadas pelo desenho do \textit{benchmark}, a contribuição do
RAG é clara na comparação qualitativa. Sem recuperação, o modelo ou inventa conteúdo
ou se recusa a responder de forma genérica; com RAG, ou cita o material real ou
admite honestamente a ausência do dado no contexto. Por exemplo, à pergunta sobre os
tópicos de um arquivo específico de um docente, o \textit{baseline} respondeu com uma
recusa genérica (``não tenho acesso a documentos específicos''), enquanto o RAG
respondeu ``não encontrei essa informação no contexto fornecido'', comportamento
factualmente correto e verificável. O trabalho aponta como próximos passos o
alinhamento dos metadados enviados ao Ragas, o saneamento do \textit{benchmark} e a
adoção de recuperação hierárquica (\textit{parent document retriever}).
